\chapter{Analisi ammortizzata}\label{ch:ammortizzata}

Quando si studia una struttura dati si è soliti calcolare il costo di ogni singola operazione nel caso peggiore, e da quello ricavare il costo di una sequenza di operazioni. Il problema di questo approccio è che molte strutture dati hanno operazioni che sono occasionalmente molto costose ma che, proprio per questo, lasciano la struttura in uno stato in cui le operazioni successive costano poco: l'analisi del caso peggiore della singola operazione non tiene conto di questo effetto e finisce per fornire stime troppo pessimistiche. L'\textbf{analisi ammortizzata} (amortized analysis) studia invece il costo di un'intera sequenza di operazioni e lo ripartisce in qualche modo tra le operazioni che la compongono. A differenza dell'analisi del caso medio non interviene alcuna probabilità, e si ottiene comunque una garanzia sul caso peggiore. In questo capitolo vediamo i tre metodi principali dell'analisi ammortizzata (aggregazione, accantonamenti e potenziale), applicandoli a due esempi: uno stack con l'operazione aggiuntiva \MultiPop{} e un contatore binario con l'operazione \Increment{}. Infine li usiamo per studiare un'applicazione reale, le \emph{tabelle dinamiche}, ovvero tabelle che vengono riallocate quando si riempiono o si svuotano troppo.

\section{L'idea di fondo}\label{sec:idea}
Consideriamo una struttura dati su cui si esegue una sequenza di $n$ operazioni. Ciò che ci interessa non è il costo della singola operazione, ma il tempo complessivo della sequenza, che indichiamo con $T(n)$. Questo tempo viene poi ripartito in qualche modo tra le $n$ operazioni: la scelta più semplice è dividerlo in parti uguali.

\begin{definizione}[Costo ammortizzato per operazione]\label{def:ammortizzato}
Sia $T(n)$ il costo totale, nel caso peggiore, di una sequenza di $n$ operazioni su una struttura dati. Il \textbf{costo ammortizzato per operazione} è
\[
\frac{T(n)}{n}.
\]
\end{definizione}

Si tratta quindi del costo che ogni operazione ha \emph{in media} all'interno della sequenza, ovvero di una media sul caso peggiore: a ogni operazione si attribuisce una parte del costo totale, anche se in quel momento la sua esecuzione è stata più economica o più costosa.

\begin{osservazione}[Non è un'analisi probabilistica]
La stima ottenuta \textbf{non è probabilistica}: non si fanno ipotesi sulla distribuzione delle operazioni e non si calcolano valori attesi. Si considera la sequenza peggiore possibile e si garantisce che, anche in quel caso, il costo medio per operazione resta contenuto, ovvero che il costo ammortizzato è un limite valido per qualsiasi sequenza di operazioni.
\end{osservazione}

\paragraph{L'approccio naive.}
Un modo naturale per calcolare la complessità di una sequenza di $n$ operazioni consiste nel calcolare il caso peggiore di una \emph{singola} operazione e nel moltiplicarlo per $n$. Vedremo che questo approccio dà luogo a un'analisi \textbf{troppo pessimistica}. Infatti, in molte strutture dati è impossibile che tutte le operazioni di una sequenza ricadano nel caso peggiore: anzi, spesso solo una parte molto limitata di esse può effettivamente incorrere in tale situazione. L'analisi ammortizzata serve proprio a tenere conto di questo fatto.

\paragraph{I tre metodi.}
Per ottenere stime più strette si usano tre tecniche:
\begin{itemize}
    \item \textbf{Metodo dell'aggregazione}: si stima il costo totale $T(n)$ della sequenza e lo si divide in parti uguali tra le $n$ operazioni (sezione~\ref{sec:aggregazione}).
    \item \textbf{Metodo degli accantonamenti} (accounting method): si assegna a ciascuna operazione un costo ammortizzato scelto da noi, in modo che la somma dei costi ammortizzati non sia mai inferiore al costo reale; l'eccedenza viene immagazzinata come \emph{credito} su singoli elementi della struttura dati (sezione~\ref{sec:accantonamenti}).
    \item \textbf{Metodo del potenziale} (potential method): come il precedente, ma il credito viene associato alla struttura dati nel suo insieme tramite una \emph{funzione potenziale} (sezione~\ref{sec:potenziale}).
\end{itemize}

\section{Due esempi di riferimento}\label{sec:esempi}
Per confrontare i tre metodi useremo sempre gli stessi due esempi, così da vedere come lo stesso problema venga affrontato in modi diversi.

\subsection{Stack con \MultiPop}\label{sec:stack}
Uno stack $S$ è una struttura dati che si accede solo dalla cima, e che supporta le seguenti operazioni elementari:
\begin{itemize}
    \item $\Push(S,x)$: inserisce l'elemento $x$ in cima allo stack.
    \item $\Pop(S)$: rimuove l'elemento in cima allo stack, se c'è; su uno stack vuoto non ha alcun effetto.
    \item $\StackEmpty(S)$: restituisce vero se lo stack è vuoto, falso altrimenti.
\end{itemize}
Ognuna di queste operazioni costa $O(1)$, e per semplicità attribuiamo costo $1$ a \Push{} e a \Pop{}. Più precisamente, contiamo soltanto i \Push{} e i \Pop{} \emph{effettivi}, ovvero quelli che inseriscono o rimuovono davvero un elemento: un \Pop{} eseguito su uno stack vuoto ha quindi costo $0$. Aggiungiamo ora l'operazione \textbf{\MultiPop}$(S,k)$, che rimuove i $k$ elementi in cima allo stack, oppure l'intero stack se questo ne contiene meno di $k$. Si tratta di un'operazione \emph{derivata}, ovvero ottenuta ripetendo più volte un'operazione elementare.

\begin{algorithm}[htbp]
\caption{L'operazione $\MultiPop(S,k)$ ripete \Pop{} finché lo stack non si svuota o finché non sono stati rimossi $k$ elementi.}\label{alg:multipop}
\begin{algorithmic}[1]
\Procedure{MultiPop}{$S,k$}
  \While{\textbf{not} \StackEmpty($S$) \textbf{and} $k>0$}
    \State \Call{Pop}{$S$}
    \State $k \gets k-1$
  \EndWhile
\EndProcedure
\end{algorithmic}
\end{algorithm}

Il costo di $\MultiPop(S,k)$ è pari al numero di iterazioni del ciclo, ovvero al numero di \Pop{} effettivamente eseguiti:
\[
\text{costo}\big(\MultiPop(S,k)\big) = \min(|S|,\,k)
\]
dove $|S|$ è il numero di elementi presenti nello stack al momento della chiamata. La figura~\ref{fig:multipop} mostra un esempio.

\begin{figure}[htbp]
\centering
\begin{widget}{stack}[items={4,8,15,16,23,42}, ops={multipop 4, multipop 7}]
\begin{tikzpicture}[>=Stealth, el/.style={draw=Accent, fill=Accent!6, minimum width=11mm, minimum height=5.5mm, inner sep=0pt, font=\small}]
  % (a) stack iniziale
  \foreach \v [count=\i from 0] in {4,8,15,16,23,42}{\node[el] at (0,{0.55*\i}) {\v};}
  \draw[thick] (-0.7,-0.3) -- (0.7,-0.3);
  \node[font=\scriptsize, TextGray, anchor=west] at (0.75,2.75) {$\leftarrow$ cima};
  \node[font=\small] at (0,-0.85) {(a)};
  % (b) dopo MULTIPOP(S,4)
  \foreach \v [count=\i from 0] in {4,8}{\node[el] at (5.2,{0.55*\i}) {\v};}
  \draw[thick] (4.5,-0.3) -- (5.9,-0.3);
  \node[font=\scriptsize, TextGray, anchor=west] at (5.95,0.55) {$\leftarrow$ cima};
  \node[font=\small] at (5.2,-0.85) {(b)};
  % (c) dopo MULTIPOP(S,7)
  \draw[thick] (9.5,-0.3) -- (10.9,-0.3);
  \node[font=\scriptsize, TextGray] at (10.2,0.2) {vuoto};
  \node[font=\small] at (10.2,-0.85) {(c)};
  % frecce
  \draw[->, thick] (1.5,1.4) -- node[above, font=\scriptsize] {$\MultiPop(S,4)$} (3.9,1.4);
  \draw[->, thick] (6.7,0.5) -- node[above, font=\scriptsize] {$\MultiPop(S,7)$} (9.1,0.5);
\end{tikzpicture}
\end{widget}
\caption{Effetto di \MultiPop{} su uno stack contenente sei elementi (a, la cima è l'elemento 42). $\MultiPop(S,4)$ rimuove i quattro elementi in cima e lascia 4 e 8 (b); la chiamata successiva $\MultiPop(S,7)$ svuota lo stack (c), perché ne erano rimasti solo due: in questo caso il costo è $\min(2,7)=2$ e non $7$.}
\label{fig:multipop}
\end{figure}

\paragraph{Analisi naive.}
Consideriamo una sequenza di $n$ operazioni (\Push, \Pop{} e \MultiPop) a partire da uno stack inizialmente vuoto. In questa situazione:
\begin{itemize}
    \item Lo stack contiene al più $n$ elementi, quindi $|S| = O(n)$.
    \item Una singola operazione costa quindi $O(n)$ nel caso peggiore, perché un \MultiPop{} può dover rimuovere $O(n)$ elementi.
    \item Il costo di $n$ operazioni è allora $n \cdot O(n) = O(n^2)$.
\end{itemize}
Questo limite superiore è corretto, ma non è \textbf{stretto}\footnote{Ovvero non è anche un limite $\Omega(n^2)$: nessuna sequenza di $n$ operazioni costa davvero quanto $n^2$.}: abbiamo ipotizzato che ogni operazione sia un \MultiPop{} costoso, e vedremo che ciò è impossibile.

\subsection{Contatore binario con \Increment}\label{sec:contatore}
Il secondo esempio è un contatore binario che conta in avanti a partire da $0$. Sia $A[0\dots k-1]$ un array di $k$ bit, in cui $A[0]$ è il bit meno significativo. Il valore memorizzato nel contatore è
\[
\mathrm{value}[A] = \sum_{i=0}^{k-1} A[i]\cdot 2^i.
\]
L'operazione \Increment{} somma $1$ al valore, modulo $2^k$:
\[
\mathrm{value}[\Increment(A)] \equiv \mathrm{value}[A] + 1 \pmod{2^k}.
\]
Ad esempio, con $k=7$ il contatore $1101011$ (che vale $107$) diventa $1101100$ (che vale $108$): gli ultimi due bit passano da $1$ a $0$, il terzo da $0$ a $1$, e gli altri restano invariati.

\begin{algorithm}[htbp]
\caption{L'operazione $\Increment(A)$ azzera i bit uguali a $1$ a partire dal meno significativo (propagando il riporto) e porta a $1$ il primo bit uguale a $0$ che incontra.}\label{alg:increment}
\begin{algorithmic}[1]
\Procedure{Increment}{$A$}
  \State $i \gets 0$
  \While{$i<k$ \textbf{and} $A[i]=1$}
    \State $A[i] \gets 0$
    \State $i \gets i+1$
  \EndWhile
  \If{$i<k$}
    \State $A[i] \gets 1$\label{alg:increment:set}
  \EndIf
\EndProcedure
\end{algorithmic}
\end{algorithm}

\noindent
Le operazioni elementari sono la modifica di un singolo bit, e ne distinguiamo due tipi:
\begin{itemize}
    \item \textbf{\SetBit}: porta un bit da $0$ a $1$.
    \item \textbf{\ResetBit}: porta un bit da $1$ a $0$.
\end{itemize}
Ciascuna ha costo $1$, quindi il costo di un \Increment{} è il numero di bit che cambiano. Notiamo che ogni \Increment{} esegue \emph{al più un} \SetBit{} (riga~\ref{alg:increment:set} dell'algoritmo~\ref{alg:increment}), e può non eseguirne nessuno solo quando il contatore passa da $1\dots1$ a $0\dots0$.

\paragraph{Analisi naive.}
Nel caso peggiore l'array contiene tutti $1$: l'incremento azzera tutti i $k$ bit, quindi
\[
\text{costo di un incremento} = O(k), \qquad \text{costo di } n \text{ incrementi} = n \cdot O(k) = O(nk).
\]
Anche in questo caso si tratta di un'analisi molto sovrabbondante: un incremento modifica tutti i $k$ bit soltanto nel passaggio da $01\dots1$ a $10\dots0$ e in quello da $1\dots1$ a $0\dots0$, cioè due volte ogni $2^k$ incrementi. Subito dopo il contatore ha al più un bit uguale a $1$, e gli incrementi successivi costano pochissimo.

\section{Metodo dell'aggregazione}\label{sec:aggregazione}
Il problema delle analisi naive viste sopra è che calcolano il caso peggiore di una singola operazione e poi lo moltiplicano per $n$. L'idea dell'aggregazione è non prendere questa scorciatoia: si stima direttamente il costo totale della sequenza, tenendo conto di come le operazioni interagiscono tra loro.

\begin{definizione}[Metodo dell'aggregazione]\label{def:aggregazione}
Il \textbf{metodo dell'aggregazione} (aggregate analysis) consiste nello stimare il costo $T(n)$ di una sequenza di $n$ operazioni e nel distribuirlo in modo equo tra le $n$ operazioni, ottenendo un costo ammortizzato
\[
\frac{T(n)}{n}
\]
per ciascuna operazione, indipendentemente dal suo tipo.
\end{definizione}

\subsection{Stack con \MultiPop}
Partiamo da uno stack inizialmente vuoto. Le operazioni \Push{} e \Pop{} sono le operazioni \emph{elementari}, mentre \MultiPop{} è un'operazione \emph{derivata}: ogni sequenza di operazioni $\langle op_1,\dots,op_n\rangle$ con $op_i \in \{\Push,\Pop,\MultiPop\}$ può essere \textbf{srotolata} in una sequenza equivalente $\langle op'_1,\dots,op'_m\rangle$ di sole operazioni elementari, con $op'_j \in \{\Push,\Pop\}$, sostituendo
\begin{center}
\begin{tabular}{rcl}
$\Pop(S)$ & $\longrightarrow$ & $\Pop(S)$ \\
$\Push(S,x)$ & $\longrightarrow$ & $\Push(S,x)$ \\
$\MultiPop(S,k)$ & $\longrightarrow$ & $\underbrace{\Pop(S),\dots,\Pop(S)}_{\min(|S|,\,k)}$
\end{tabular}
\end{center}
Poiché ogni operazione elementare costa $1$, le due sequenze hanno lo stesso costo, ovvero il numero $m$ di operazioni elementari:
\[
\text{costo}\big(\langle op_1,\dots,op_n\rangle\big) = \text{costo}\big(\langle op'_1,\dots,op'_m\rangle\big) = m.
\]

\paragraph{Esempio.}
Consideriamo la sequenza $\Push(a)$, $\Push(b)$, $\Push(c)$, $\MultiPop(4)$, $\Push(d)$, $\Pop$, $\dots$ e osserviamo come viene srotolata in figura~\ref{fig:srotolamento}. Quando viene eseguito $\MultiPop(4)$ lo stack contiene soltanto tre elementi, quindi l'operazione si traduce in $\min(3,4)=3$ operazioni di \Pop.

\begin{figure}[htbp]
\centering
\begin{tikzpicture}[>=Stealth, every node/.style={font=\scriptsize\sffamily}, opl/.style={font=\scriptsize, TextGray}]
  % operazioni originali
  \node (t1) at (0,2.2) {PUSH(a)};
  \node (t2) at (2,2.2) {PUSH(b)};
  \node (t3) at (4,2.2) {PUSH(c)};
  \node (t4) at (6.6,2.2) {MULTIPOP(4)};
  \node (t5) at (9.4,2.2) {PUSH(d)};
  \node (t6) at (11.4,2.2) {POP()};
  \foreach \i/\x in {1/0,2/2,3/4,4/6.6,5/9.4,6/11.4}{\node[opl] at (\x,2.75) {$op_\i$};}
  % operazioni elementari
  \node (b1) at (0,0) {PUSH(a)};
  \node (b2) at (2,0) {PUSH(b)};
  \node (b3) at (4,0) {PUSH(c)};
  \node (b4) at (5.5,0) {POP()};
  \node (b5) at (6.6,0) {POP()};
  \node (b6) at (7.7,0) {POP()};
  \node (b7) at (9.4,0) {PUSH(d)};
  \node (b8) at (11.4,0) {POP()};
  \foreach \i/\x in {1/0,2/2,3/4,4/5.5,5/6.6,6/7.7,7/9.4,8/11.4}{\node[opl] at (\x,-0.5) {$op'_\i$};}
  % frecce
  \draw[->, Accent] (t1.south) -- (b1.north);
  \draw[->, Accent] (t2.south) -- (b2.north);
  \draw[->, Accent] (t3.south) -- (b3.north);
  \draw[->, Accent] (t4.south) -- (b4.north);
  \draw[->, Accent] (t4.south) -- (b5.north);
  \draw[->, Accent] (t4.south) -- (b6.north);
  \draw[->, Accent] (t5.south) -- (b7.north);
  \draw[->, Accent] (t6.south) -- (b8.north);
  \draw[decorate, decoration={brace, mirror, amplitude=4pt}] (5.0,-0.85) -- (8.2,-0.85) node[midway, below=4pt] {$\min(|S|,4)=3$};
\end{tikzpicture}
\caption{Una sequenza di operazioni sullo stack (in alto) e la corrispondente sequenza di sole operazioni elementari (in basso): $\MultiPop(4)$, eseguito su uno stack che contiene tre elementi, si srotola in tre \Pop{} (parentesi graffa). La sequenza originale ha $6$ operazioni e quella srotolata ne ha $8$; ognuna delle $8$ costa $1$.}
\label{fig:srotolamento}
\end{figure}

\paragraph{Un secondo esempio.}
Prendiamo la sequenza $\Push(1)$, $\Push(2)$, $\Push(3)$, $\Pop$, $\Push(4)$, $\MultiPop(2)$. Le prime cinque operazioni costano $1$ ciascuna, ovvero $4$ \Push{} e $1$ \Pop, mentre il \MultiPop{} rimuove $4$ e $2$ e costa quindi $2$: il costo totale è $4+1+2=7$. Se aggiungiamo un altro $\MultiPop(2)$, lo stack contiene solo l'elemento $1$, quindi si traduce in un solo \Pop{} e il costo totale sale a $8$. Notiamo che un \MultiPop{} costa tanto quanti sono gli elementi che toglie, e che non può togliere più elementi di quelli inseriti.

\paragraph{Analisi.}
Questa osservazione è il punto chiave: ogni elemento può essere rimosso al più una volta per ogni volta che è stato inserito. Nella sequenza srotolata, quindi, il numero di \Pop{} non supera il numero di \Push, e il numero di \Push{} della sequenza srotolata coincide con quello della sequenza originale, a sua volta al più $n$:
\[
\#\Pop\big(\langle op'_1,\dots,op'_m\rangle\big) \le \#\Push\big(\langle op'_1,\dots,op'_m\rangle\big) = \#\Push\big(\langle op_1,\dots,op_n\rangle\big) \le n.
\]
Detta $\langle op_1,\dots,op_n\rangle$ la sequenza più costosa, si ottiene
\[
T(n) = \text{costo}\big(\langle op'_1,\dots,op'_m\rangle\big) = \#\Pop + \#\Push \le n + n = 2n
\]
e quindi
\[
\text{costo ammortizzato per operazione} = \frac{T(n)}{n} \le \frac{2n}{n} = 2 = O(1).
\]
Lo stesso costo ammortizzato vale per \Push, \Pop{} e \MultiPop{}, anche se un singolo \MultiPop{} può costare molto: è il totale della sequenza che resta limitato da $2n$, non la singola operazione. Rispetto all'analisi naive, che dava $O(n^2)$, il costo di $n$ operazioni scende a $O(n)$.

\begin{esercizio}[Sequenza più costosa sullo stack]
Dimostrare che la sequenza più costosa di $n$ operazioni \Push, \Pop{} e \MultiPop, a partire da uno stack vuoto, ha costo
\[
T(n) \le \begin{cases} 1 & \text{se } n = 1 \\ 2(n-1) & \text{se } n \ge 2. \end{cases}
\]
\end{esercizio}

\paragraph{Soluzione.}
Se $n=1$ le sequenze possibili sono tre: $\langle\Push\rangle$, di costo $1$, $\langle\Pop\rangle$, di costo $0$ perché lo stack è vuoto, e $\langle\MultiPop(k)\rangle$ con $k = 1,2,\dots$, anch'esse di costo $0$. Quindi $T(1) \le 1$.

Sia ora $n \ge 2$. Indichiamo con $p$ il numero di \Push{} della sequenza e con $r$ il numero di elementi rimossi da \Pop{} e \MultiPop. Ogni elemento può essere rimosso solo dopo essere stato inserito, quindi $r \le p$ e
\[
T(n) = p + r \le 2p.
\]
Distinguiamo due casi:
\begin{itemize}
    \item Se $r \ge 1$, almeno una delle operazioni è un \Pop{} o un \MultiPop, quindi $p \le n-1$ e $T(n) \le 2(n-1)$.
    \item Se $r = 0$, il costo è dato dai soli \Push, e poiché $n - 2 \ge 0$ si ha $T(n) = p \le n \le n + (n-2) = 2(n-1)$.
\end{itemize}
Il limite $2(n-1)$ si raggiunge davvero: basta eseguire $n-1$ operazioni \Push{} seguite da $\MultiPop(S,n-1)$, che costano $(n-1)+(n-1)=2(n-1)$.

\subsection{Contatore binario con \Increment}\label{sec:contatore-aggr}
Supponiamo ora che il contatore sia inizialmente nullo, e studiamo quanti bit cambiano in una sequenza di $n$ incrementi. Il punto è che non tutti i bit cambiano a ogni chiamata, come si vede dalla figura~\ref{fig:contatore}:
\begin{itemize}
    \item $A[0]$ cambia a ogni incremento, ovvero $n$ volte.
    \item $A[1]$ cambia solo una volta ogni due incrementi, ovvero $\lfloor n/2\rfloor$ volte.
    \item $A[2]$ cambia solo una volta ogni quattro incrementi, ovvero $\lfloor n/2^2\rfloor$ volte.
\end{itemize}
In generale, per $i = 0,1,\dots,k-1$ il bit $A[i]$ cambia $\lfloor n/2^i \rfloor$ volte nell'arco di $n$ incrementi\footnote{Il bit $A[i]$ cambia nel passaggio dal valore $v$ al valore $v+1$ se e solo se $v+1$ è un multiplo di $2^i$. Partendo da $0$, i valori $v+1$ percorsi sono $1,\dots,n$, tra i quali i multipli di $2^i$ sono esattamente $\lfloor n/2^i\rfloor$.}; per $i \ge k$ il bit non esiste e non può cambiare.

\begin{figure}[htbp]
\centering
\begin{widget}{binary-counter}[bits=8, steps=16]
\begin{tikzpicture}[
  cb0/.style={minimum width=7mm, minimum height=4.4mm, inner sep=0pt, font=\scriptsize\ttfamily, draw=RuleGray, fill=white},
  cb1/.style={minimum width=7mm, minimum height=4.4mm, inner sep=0pt, font=\scriptsize\ttfamily, draw=RuleGray, fill=Accent!25},
  lab/.style={font=\scriptsize, TextGray},
  x=1mm, y=1mm]
  \newcommand{\cbrow}[3]{%
    \node[lab] at (-12,-5*#1) {#1};
    \foreach \b/\f [count=\j from 0] in {#2}{\node[cb\f] at (7*\j,-5*#1) {\b};}
    \node[lab] at (62,-5*#1) {#3};}
  % intestazione
  \foreach \j/\x in {7/0,6/7,5/14,4/21,3/28,2/35,1/42,0/49}{\node[lab] at (\x,5) {$A[\j]$};}
  \node[lab] at (-12,5) {valore};
  \node[lab, align=center] at (62,6) {costo\\totale};
\cbrow{0}{0/0,0/0,0/0,0/0,0/0,0/0,0/0,0/0}{0}
\cbrow{1}{0/0,0/0,0/0,0/0,0/0,0/0,0/0,1/1}{1}
\cbrow{2}{0/0,0/0,0/0,0/0,0/0,0/0,1/1,0/1}{3}
\cbrow{3}{0/0,0/0,0/0,0/0,0/0,0/0,1/0,1/1}{4}
\cbrow{4}{0/0,0/0,0/0,0/0,0/0,1/1,0/1,0/1}{7}
\cbrow{5}{0/0,0/0,0/0,0/0,0/0,1/0,0/0,1/1}{8}
\cbrow{6}{0/0,0/0,0/0,0/0,0/0,1/0,1/1,0/1}{10}
\cbrow{7}{0/0,0/0,0/0,0/0,0/0,1/0,1/0,1/1}{11}
\cbrow{8}{0/0,0/0,0/0,0/0,1/1,0/1,0/1,0/1}{15}
\cbrow{9}{0/0,0/0,0/0,0/0,1/0,0/0,0/0,1/1}{16}
\cbrow{10}{0/0,0/0,0/0,0/0,1/0,0/0,1/1,0/1}{18}
\cbrow{11}{0/0,0/0,0/0,0/0,1/0,0/0,1/0,1/1}{19}
\cbrow{12}{0/0,0/0,0/0,0/0,1/0,1/1,0/1,0/1}{22}
\cbrow{13}{0/0,0/0,0/0,0/0,1/0,1/0,0/0,1/1}{23}
\cbrow{14}{0/0,0/0,0/0,0/0,1/0,1/0,1/1,0/1}{25}
\cbrow{15}{0/0,0/0,0/0,0/0,1/0,1/0,1/0,1/1}{26}
\cbrow{16}{0/0,0/0,0/0,1/1,0/1,0/1,0/1,0/1}{31}
\end{tikzpicture}
\end{widget}
\caption{Un contatore a $8$ bit durante $16$ incrementi consecutivi, a partire da $0$. A sinistra il valore del contatore, al centro i bit $A[7],\dots,A[0]$ (i bit evidenziati sono quelli che sono cambiati nel passaggio dal valore precedente), a destra il costo totale, ovvero il numero complessivo di bit modificati dall'inizio. $A[0]$ cambia a ogni incremento, $A[1]$ una volta ogni due, $A[2]$ una volta ogni quattro e così via; il costo totale è sempre minore del doppio del numero di incrementi.}
\label{fig:contatore}
\end{figure}

Il costo totale è il numero complessivo di bit che cambiano, quindi
\[
T(n) = \sum_{i=0}^{k-1} \left\lfloor \frac{n}{2^i} \right\rfloor \le \sum_{i=0}^{k-1} \frac{n}{2^i} < \sum_{i=0}^{\infty} \frac{n}{2^i} = n \sum_{i=0}^{\infty} \frac{1}{2^i} = 2n
\]
dove abbiamo usato il fatto che la serie geometrica $\sum_{i\ge 0} 2^{-i}$ vale $2$. Ne segue che
\[
\text{costo ammortizzato per operazione} = \frac{T(n)}{n} < \frac{2n}{n} = 2 = O(1)
\]
contro il $O(k)$ dell'analisi naive. In figura, ad esempio, $16$ incrementi costano $31 < 32 = 2\cdot 16$.

\begin{esercizio}[Costo di una sequenza di incrementi]
Dimostrare che la sequenza di $n$ incrementi su un contatore binario a $k$ bit, inizialmente nullo, ha costo $T(n)$ tale che
\[
n \le T(n) \le 2n-1
\]
dove le due disuguaglianze non sono in generale rafforzabili.
\end{esercizio}

\paragraph{Soluzione.}
Ogni incremento modifica almeno un bit, quindi $T(n) \ge n$. Per il limite superiore basta il calcolo appena fatto con l'aggregazione, che dà $T(n) < 2n$: poiché $T(n)$ è un numero intero, ne segue
\[
T(n) \le 2n - 1.
\]
Le due disuguaglianze non si possono migliorare. Con $k=1$ ogni incremento modifica esattamente un bit, quindi $T(n)=n$. Per la seconda consideriamo $n = 2^i$ con $1 \le i \le k-1$: i bit $A[j]$ con $j > i$ non cambiano mai, perché $\lfloor 2^i/2^j \rfloor = 0$, e quindi
\[
T(2^i) = 2^i + \left\lfloor \frac{2^i}{2} \right\rfloor + \dots + \left\lfloor \frac{2^i}{2^{i-1}} \right\rfloor + \left\lfloor \frac{2^i}{2^i} \right\rfloor = 2^i + 2^{i-1} + \dots + 2 + 1 = 2^{i+1} - 1 = 2n - 1.
\]
Ad esempio $T(4) = 1+2+1+3 = 7 = 2\cdot 4 -1$.

\subsection{Contatore non necessariamente nullo}
Se il contatore parte da una configurazione qualsiasi, il bit $A[i]$ continua a cambiare una volta ogni $2^i$ incrementi, ma il primo cambiamento può arrivare subito. Per la stessa ragione della nota precedente, il numero di volte in cui $A[i]$ cambia è il numero di multipli di $2^i$ tra $n$ interi consecutivi, che è al più $\lceil n/2^i \rceil$. Si ottiene quindi
\[
T(n) \le \sum_{i=0}^{k-1} \left\lceil \frac{n}{2^i} \right\rceil \le \sum_{i=0}^{k-1} \left(\frac{n}{2^i} + 1\right) = n \sum_{i=0}^{k-1} \frac{1}{2^i} + k < 2n + k.
\]
Il costo di $n$ incrementi è quindi $O(n+k)$: se $n = \Omega(k)$, ovvero se si eseguono almeno tanti incrementi quanti sono i bit, il costo è $O(n)$ qualunque sia la configurazione di partenza, e il costo ammortizzato resta $O(1)$.

\paragraph{Esempio.}
Partiamo dal contatore a $6$ bit $011011$ ed eseguiamo $n=5$ incrementi, che producono nell'ordine $011100$, $011101$, $011110$, $011111$ e $100000$. Il bit $A[0]$ cambia $5$ volte, $A[1]$ cambia $3$ volte (nel primo, nel terzo e nel quinto incremento), $A[2]$ due volte, mentre $A[3]$, $A[4]$ e $A[5]$ cambiano una volta soltanto, nell'ultimo incremento. Questi numeri coincidono con i limiti $\lceil 5/2^i \rceil = 5,3,2,1,1,1$, e il costo totale è $13 < 2n + k = 16$.

\begin{esercizio}[\MultiPush]
Se l'insieme delle operazioni sullo stack includesse un'operazione $\MultiPush(S,x_1,\dots,x_k)$, che inserisce $k$ elementi nello stack, il limite $O(1)$ sul costo ammortizzato delle operazioni sullo stack sarebbe ancora valido?
\end{esercizio}

\paragraph{Soluzione.}
No, in generale. Un \MultiPush{} costa $k$ e può inserire fino a $k$ elementi, quindi una sequenza di $n$ operazioni \MultiPush{} su uno stack vuoto costa $nk$ e ha costo ammortizzato $k$, che non è $O(1)$ se $k$ non è limitato da una costante. L'argomento del paragrafo precedente cade perché non è più vero che il numero di \Push{} è al più $n$: la sequenza ne esegue anche $nk$.

\begin{esercizio}[\Decrement]
Si dimostri che, se un'operazione \Decrement{} fosse inclusa nell'esempio del contatore a $k$ bit, $n$ operazioni potrebbero richiedere tempo $\Theta(nk)$.
\end{esercizio}

\paragraph{Soluzione.}
Ogni operazione modifica al più $k$ bit, quindi $n$ operazioni costano $O(nk)$. Per il limite inferiore partiamo da $0\dots0$ ed eseguiamo alternativamente \Decrement{} e \Increment: il primo porta il contatore a $1\dots1$ (modulo $2^k$) modificando $k$ bit, il secondo lo riporta a $0\dots0$ modificandone ancora $k$. Ogni operazione costa quindi $k$ e il totale è $nk = \Theta(nk)$. L'analisi dell'incremento non si estende perché sfrutta il fatto che un bit può essere azzerato solo dopo essere stato portato a $1$; con \Decrement{} il contatore può tornare indietro e il caso peggiore si ripete a ogni passo.

\begin{esercizio}[Costo pari a una potenza di $2$]\label{es:potenze}
Una sequenza di $n$ operazioni viene eseguita su una struttura dati. L'$i$-esima operazione costa $i$ se $i$ è una potenza esatta di $2$, altrimenti costa $1$. Applicare il metodo dell'aggregazione per determinare il costo ammortizzato per operazione. Lo stesso esercizio verrà risolto con il metodo degli accantonamenti (sezione~\ref{sec:accantonamenti}) e con quello del potenziale (sezione~\ref{sec:potenziale}).
\end{esercizio}

\paragraph{Soluzione.}
Indichiamo con $c_i$ il costo dell'$i$-esima operazione, quindi $c_i = i$ se $i = 2^j$ per un certo $j \in \mathbb{N}$ e $c_i = 1$ altrimenti. Separiamo nella somma le operazioni che costano $1$ da quelle che costano $i$:
\[
T(n) = \sum_{i=1}^{n} c_i = \sum_{\substack{i=1 \\ i \notin 2^{\mathbb{N}}}}^{n} 1 + \sum_{\substack{i=1 \\ i \in 2^{\mathbb{N}}}}^{n} i.
\]
Tra $1$ e $n$ ci sono $\lfloor \lg n \rfloor + 1$ potenze di $2$, ovvero $2^0, 2^1,\dots,2^{\lfloor \lg n\rfloor}$\footnote{Con $\lg$ si indica il logaritmo in base $2$.}, quindi la prima somma vale $n - (\lfloor \lg n \rfloor + 1)$, mentre la seconda è una somma geometrica:
\[
\sum_{\substack{i=1 \\ i \in 2^{\mathbb{N}}}}^{n} i = \sum_{j=0}^{\lfloor \lg n \rfloor} 2^j = 2^{\lfloor \lg n \rfloor + 1} - 1.
\]
Si ottiene allora
\[
T(n) = n - (\lfloor \lg n \rfloor + 1) + 2^{\lfloor \lg n \rfloor + 1} - 1 < n + 2^{\lfloor \lg n \rfloor + 1} \le n + 2n = 3n
\]
dove abbiamo usato $2^{\lfloor \lg n \rfloor} \le n$. Pertanto
\[
\text{costo ammortizzato per operazione} = \frac{T(n)}{n} < \frac{3n}{n} = 3.
\]

\section{Metodo degli accantonamenti}\label{sec:accantonamenti}
Il metodo dell'aggregazione assegna lo stesso costo ammortizzato a tutte le operazioni, e questo può essere restrittivo quando le operazioni sono di natura diversa. Il metodo degli accantonamenti (accounting method) permette invece di assegnare a ciascuna operazione un costo ammortizzato scelto da noi, anche diverso da operazione a operazione.

\begin{definizione}[Costi ammortizzati]\label{def:accantonamenti}
Data una sequenza di operazioni $op_1, op_2, \dots, op_n$, indichiamo con
\[
c_i = \text{costo reale}(op_i), \qquad \hat{c}_i = \text{costo ammortizzato}(op_i)
\]
dove i costi ammortizzati $\hat{c}_i$ sono \emph{definiti da noi}. L'\textbf{obiettivo} è definirli in modo tale che, per qualsiasi sequenza di operazioni, valga
\[
T(n) = \sum_{i=1}^{n} c_i \le \sum_{i=1}^{n} \hat{c}_i.
\]
\end{definizione}

In questo modo la somma dei costi ammortizzati è un limite superiore al costo reale, e basta stimare la prima per avere una stima della seconda.

\paragraph{L'idea del credito.}
Per garantire la disuguaglianza si ragiona come in un conto in banca, usando il costo ammortizzato come ``prezzo'' addebitato a ogni operazione:
\begin{itemize}
    \item Se $\hat{c}_i > c_i$, delle $\hat{c}_i$ unità di costo addebitate, $c_i$ servono a pagare il costo reale di $op_i$, mentre le restanti $\hat{c}_i - c_i$ vengono \textbf{immagazzinate} su elementi specifici della struttura dati, come credito.
    \item Se $c_i > \hat{c}_i$, la differenza $c_i - \hat{c}_i$ viene \textbf{recuperata} dai crediti immagazzinati nella struttura dati.
\end{itemize}
Perché il metodo funzioni, il credito deve essere sufficiente a coprire ogni operazione ``in perdita''. Per formalizzare questa idea, sia $D_0$ la configurazione iniziale di una struttura dati, sia $op_1,\dots,op_n$ una sequenza di operazioni e sia $D_{j}$ la configurazione che si ottiene da $D_{j-1}$ applicando $op_{j}$, ovvero
\[
D_0 \xrightarrow{op_1} D_1 \xrightarrow{op_2} D_2 \xrightarrow{op_3} \cdots \xrightarrow{op_n} D_n.
\]
Indichiamo con $B_j$ il \textbf{credito totale} accumulato dopo l'esecuzione delle prime $j$ operazioni, ovvero la differenza tra quanto è stato addebitato e quanto è stato effettivamente speso:
\[
B_j = \sum_{i=1}^{j} (\hat{c}_i - c_i), \qquad B_0 = 0.
\]

\begin{lemma}[Correttezza degli accantonamenti]\label{lem:accantonamenti}
Se i costi ammortizzati sono scelti in modo che
\[
B_j \ge 0 \qquad \text{per } j = 0, 1, \dots, n
\]
allora $\sum_{i=1}^{j} c_i \le \sum_{i=1}^{j} \hat{c}_i$ per ogni $j$, e in particolare $\sum_{i=1}^{n} c_i \le \sum_{i=1}^{n} \hat{c}_i$.
\end{lemma}

\paragraph{Dimostrazione.}
Per ogni $j = 1,\dots,n$ si ha, per ipotesi,
\[
0 \le B_j = \sum_{i=1}^{j} (\hat{c}_i - c_i) = \sum_{i=1}^{j} \hat{c}_i - \sum_{i=1}^{j} c_i, \qquad \text{quindi} \qquad \sum_{i=1}^{j} c_i \le \sum_{i=1}^{j} \hat{c}_i.
\]
Scegliendo $j = n$ si ottiene la disuguaglianza sull'intera sequenza. \hfill $\square$

In pratica, per applicare il metodo basta scegliere i costi $\hat{c}_i$ e verificare che il credito totale $B_j$ non diventi mai negativo: in questo modo la somma dei costi ammortizzati è una stima per eccesso del costo effettivo della sequenza.

\begin{osservazione}[I crediti esistono solo nell'analisi]
Il credito è uno strumento di analisi: non compare nel codice e non occorre memorizzarlo da nessuna parte (ad esempio, non serve aggiungere un campo $x.\mathrm{credit}$ a ogni elemento). Serve solo a convincerci che i costi $\hat{c}_i$ scelti sono sufficienti.
\end{osservazione}

\subsection{Stack con \MultiPop}
Assumiamo di partire da uno stack vuoto e scegliamo i costi ammortizzati
\[
\hat{c}_{\Push} = 2, \qquad \hat{c}_{\Pop} = \hat{c}_{\MultiPop} = 0
\]
a fronte dei costi reali $c_{\Push} = 1$, $c_{\Pop} = 1$ e $c_{\MultiPop} = \min(|S|,k)$. Questi costi hanno una lettura molto intuitiva: un \Push{} addebita $2$ unità, di cui $1$ paga il costo reale dell'operazione e l'altra viene \textbf{assegnata all'elemento} appena inserito, che porta quindi con sé una unità di credito. Quando l'elemento viene rimosso, con un \Pop{} o con un \MultiPop, il costo reale della rimozione si paga proprio con quell'unità, senza addebitare nulla all'operazione. Il credito è quindi $|S|$, cioè una unità per ogni elemento presente nello stack.

\begin{esempiofig}
Consideriamo uno stack con quattro elementi, ognuno dei quali porta con sé una moneta, cioè l'unità di credito accantonata dal suo \Push. Eseguiamo $\MultiPop(S,3)$: il costo reale è $3$, ma al \MultiPop{} non addebitiamo nulla. Le tre monete degli elementi rimossi bastano a pagare le tre operazioni \Pop, e sullo stack resta un solo elemento con la sua moneta: il credito è sempre uguale al numero di elementi presenti.
\tcblower
\begin{widget}{stack}[items={x,x,x,x}, ops={multipop 3}, coins]
\begin{tikzpicture}[scale=0.76, transform shape, >=Stealth,
  el/.style={draw=Accent, fill=Accent!6, minimum width=8mm, minimum height=5.5mm, inner sep=0pt, font=\small},
  coin/.style={circle, draw=ObsColor, fill=ObsColor!25, minimum size=3.6mm, inner sep=0pt, font=\tiny}]
  % stack prima
  \foreach \i in {0,1,2,3}{\node[el] at (0,{0.6*\i}) {x}; \node[coin] at (0.75,{0.6*\i}) {1};}
  \draw[thick] (-0.6,-0.32) -- (1.1,-0.32);
  \node[font=\scriptsize] at (0.25,-0.8) {prima};
  % freccia
  \draw[->, thick] (1.5,1.2) -- node[above, font=\scriptsize] {$\MultiPop(S,3)$} (3.7,1.2);
  % stack dopo
  \node[el] at (4.6,0) {x}; \node[coin] at (5.35,0) {1};
  \draw[thick] (4.0,-0.32) -- (5.7,-0.32);
  \node[font=\scriptsize] at (4.85,-0.8) {dopo};
  % monete usate
  \foreach \i in {0,1,2}{\node[coin, fill=ObsColor!10, draw=ObsColor!50] at (7.0,{0.6*\i}) {1};}
  \node[font=\scriptsize, TextGray, align=center] at (7.0,-0.8) {pagano\\le $3$ \Pop};
\end{tikzpicture}
\end{widget}
\end{esempiofig}

Verifichiamo che il credito non sia mai negativo. In ogni istante $j$ vale
\[
B_j = \sum_{i=1}^{j} \hat{c}_i - \sum_{i=1}^{j} c_i = |S_j| \ge 0
\]
dove $|S_j|$ è la dimensione dello stack dopo $j$ operazioni: un \Push{} aggiunge $2-1=1$ al primo membro e alla dimensione, mentre un \Pop{} o un \MultiPop{} che rimuovono $k'$ elementi aggiungono $0-k'=-k'$ al primo membro e tolgono $k'$ alla dimensione. Per il lemma~\ref{lem:accantonamenti} segue
\[
\sum_{i=1}^{n} c_i \le \sum_{i=1}^{n} \hat{c}_i = 2\cdot \#\Push \le 2n
\]
e quindi, come con l'aggregazione, il costo di $n$ operazioni è $O(n)$. Si noti che questa volta il costo ammortizzato \emph{non} è lo stesso per tutte le operazioni: vale $2$ per \Push{} e $0$ per le altre due, ma questo è proprio ciò che rende il metodo più flessibile.

\subsection{Contatore binario con \Increment}
Consideriamo il contatore inizialmente nullo e usiamo come costo reale il numero di bit modificati. Come costi ammortizzati delle sottooperazioni scegliamo
\[
\hat{c}_{\SetBit} = 2, \qquad \hat{c}_{\ResetBit} = 0.
\]
Un \SetBit{} addebita quindi $2$ unità: $1$ serve a pagare l'operazione e l'altra viene accantonata \textbf{sul bit stesso}. Poiché ogni bit uguale a $1$ è stato portato a $1$ da un \SetBit, ogni bit uguale a $1$ porta con sé una unità di credito, che usiamo per pagare il \ResetBit{} quando lo riportiamo a $0$. Un \Increment{} esegue al più un \SetBit{} (e un numero qualsiasi di \ResetBit, che non vengono addebitati), quindi
\[
\hat{c}_{\Increment} = 2\cdot \#\SetBit = \begin{cases} 0 & \text{se non esegue alcun } \SetBit \\ 2 & \text{se ne esegue uno} \end{cases} \;\le\; 2.
\]
Ad esempio, nel passaggio da $0111$ a $1000$ i tre bit azzerati usano le loro tre unità di credito, mentre il bit portato a $1$ paga una unità e ne conserva un'altra.

Il credito immagazzinato in ogni istante è il numero di bit uguali a $1$ presenti nel contatore, quindi non è mai negativo:
\[
B_j = \sum_{i=1}^{j} \hat{c}_i - \sum_{i=1}^{j} c_i = \#\{\text{bit del contatore uguali a } 1\} \ge 0.
\]
Per il lemma~\ref{lem:accantonamenti} vale allora
\[
\sum_{i=1}^{n} c_i \le \sum_{i=1}^{n} \hat{c}_i \le 2n
\]
cioè il costo di $n$ incrementi è $O(n)$ e il costo ammortizzato di un incremento è $O(1)$, ritrovando il risultato dell'aggregazione.

\begin{esercizio}[Costo pari a una potenza di $2$, con gli accantonamenti]
Risolvere l'esercizio~\ref{es:potenze} con il metodo degli accantonamenti.
\end{esercizio}

\paragraph{Soluzione.}
Scegliamo i costi ammortizzati
\[
\hat{c}_1 = 1, \qquad \hat{c}_i = 2 \;\text{ se } i = 2^j \text{ con } j \ge 1, \qquad \hat{c}_i = 3 \;\text{ altrimenti}.
\]
Un'operazione che costa $1$ (ovvero con $i$ non potenza di $2$) paga $1$ per sé e accantona $2$ unità di credito. Un'operazione $i = 2^j$ con $j \ge 1$ costa $2^j$, ma tra l'operazione $2^{j-1}$ e l'operazione $2^j$ ci sono $2^{j-1}-1$ operazioni da $1$, che hanno accantonato complessivamente $2(2^{j-1}-1) = 2^j - 2$ unità: queste pagano quasi tutto il costo $2^j$ e le $2$ unità addebitate coprono la parte restante. In questo modo il credito accumulato dopo $i$ operazioni è
\[
B_i = 2\left(i - 2^{\lfloor \lg i \rfloor}\right) \ge 0
\]
(ad esempio $2$ dopo l'operazione $3$, $0$ dopo la $4$, $6$ dopo la $7$ e di nuovo $0$ dopo la $8$) e non è mai negativo. Per il lemma~\ref{lem:accantonamenti}
\[
T(n) \le \sum_{i=1}^{n} \hat{c}_i \le 3n
\]
e il costo ammortizzato per operazione è al più $3$.

\section{Metodo del potenziale}\label{sec:potenziale}
Il metodo degli accantonamenti associa il credito a singoli elementi della struttura dati. Il metodo del potenziale (potential method) segue la stessa logica, ma immagazzina il credito nella struttura dati \emph{nel suo insieme}, come se fosse un'energia potenziale che può essere rilasciata per pagare le operazioni future.

\begin{definizione}[Funzione potenziale]\label{def:potenziale}
Una \textbf{funzione potenziale} associa a ogni configurazione $D$ della struttura dati un numero reale
\[
D \longmapsto \Phi(D) \in \mathbb{R}
\]
detto \textbf{potenziale} di $D$. Data la sequenza di configurazioni $D_0 \xrightarrow{op_1} D_1 \xrightarrow{op_2} \cdots \xrightarrow{op_n} D_n$ e detto $c_i$ il costo reale di $op_i$, il \textbf{costo ammortizzato} di $op_i$ rispetto a $\Phi$ è
\[
\hat{c}_i = c_i + \Phi(D_i) - \Phi(D_{i-1}) = c_i + \Delta\Phi_i
\]
dove $\Delta\Phi_i = \Phi(D_i) - \Phi(D_{i-1})$ è la \textbf{variazione di potenziale} causata da $op_i$.
\end{definizione}

\begin{figure}[htbp]
\centering
\begin{tikzpicture}[>=Stealth, st/.style={draw=Accent, fill=Accent!6, rounded corners=2pt, minimum width=10mm, minimum height=7mm, font=\small}]
  \node[st] (d0) at (0,0) {$D_0$};
  \node[st] (d1) at (3,0) {$D_1$};
  \node[st] (d2) at (6,0) {$D_2$};
  \node (dd) at (8.6,0) {$\cdots$};
  \node[st] (dn) at (11,0) {$D_n$};
  \draw[->] (d0) -- node[above, font=\small] {$op_1$} (d1);
  \draw[->] (d1) -- node[above, font=\small] {$op_2$} (d2);
  \draw[->] (d2) -- node[above, font=\small] {$op_3$} (dd);
  \draw[->] (dd) -- node[above, font=\small] {$op_n$} (dn);
  \foreach \i/\x in {0/0,1/3,2/6,n/11}{
    \draw[FigColor, dotted] (\x,-0.4) -- (\x,-0.95);
    \node[FigColor, font=\small] at (\x,-1.3) {$\Phi(D_\i)$};}
\end{tikzpicture}
\caption{Il metodo del potenziale associa a ogni configurazione $D_i$ della struttura dati un valore $\Phi(D_i)$: il costo ammortizzato dell'operazione $op_i$ è il suo costo reale più la differenza tra il potenziale dopo e prima dell'operazione.}
\label{fig:potenziale}
\end{figure}

Intuitivamente, se $\Delta\Phi_i > 0$ l'operazione viene sovrapagata e il potenziale della struttura dati aumenta, mentre se $\Delta\Phi_i < 0$ l'operazione viene sottopagata e la diminuzione del potenziale paga la parte di costo reale che non è coperta dal costo ammortizzato. Sommando i costi ammortizzati, i termini di potenziale si semplificano a coppie:
\begin{align*}
\sum_{i=1}^{n} \hat{c}_i &= \sum_{i=1}^{n} \big[c_i + \Phi(D_i) - \Phi(D_{i-1})\big] \\
&= \big(c_1 + \Phi(D_1) - \Phi(D_0)\big) + \big(c_2 + \Phi(D_2) - \Phi(D_1)\big) + \cdots + \big(c_n + \Phi(D_n) - \Phi(D_{n-1})\big)
\end{align*}
dove ogni $\Phi(D_i)$, con $0 < i < n$, compare una volta con segno positivo e una con segno negativo (somma telescopica), da cui
\begin{equation}\label{eq:telescopica}
\sum_{i=1}^{n} \hat{c}_i = \sum_{i=1}^{n} c_i + \Phi(D_n) - \Phi(D_0).
\end{equation}

\begin{lemma}[Correttezza del potenziale]\label{lem:potenziale}
Vale $\displaystyle\sum_{i=1}^{n} \hat{c}_i \ge \sum_{i=1}^{n} c_i$ se e solo se $\Phi(D_n) \ge \Phi(D_0)$.
\end{lemma}

\paragraph{Dimostrazione.}
Per la~\eqref{eq:telescopica}, la disuguaglianza $\sum \hat{c}_i \ge \sum c_i$ equivale a $\sum c_i + \Phi(D_n) - \Phi(D_0) \ge \sum c_i$, ovvero a $\Phi(D_n) - \Phi(D_0) \ge 0$. \hfill $\square$

Nella pratica non si sa in anticipo quante operazioni verranno eseguite, quindi si richiede la condizione per ogni $k$:
\[
\Phi(D_k) \ge \Phi(D_0) \qquad \text{per } k = 1,2,\dots,n.
\]
Di solito si definisce $\Phi(D_0) = 0$ e si dimostra che $\Phi(D_k) \ge 0$ per ogni $k$. Notiamo che i costi ammortizzati dipendono dalla funzione potenziale scelta: funzioni diverse danno costi diversi, tutti però validi come limiti superiori al costo reale. La scelta migliore dipende dai limiti che si vogliono dimostrare.

\subsection{Stack con \MultiPop}
Partendo da uno stack inizialmente vuoto, definiamo il potenziale come il numero di elementi dello stack:
\[
\Phi(S) = |S|.
\]
Se $S_0$ è lo stack vuoto si ha $\Phi(S_0) = 0$, e poiché la dimensione di uno stack non è mai negativa vale $\Phi(S) \ge \Phi(S_0)$ per ogni configurazione. Calcoliamo i costi ammortizzati delle tre operazioni, indicando con $S_{i-1}$ e $S_i$ lo stack prima e dopo l'operazione:
\begin{align*}
\hat{c}_{\Push} &= c_{\Push} + |S_i| - |S_{i-1}| = 1 + 1 = 2 && \text{perché } |S_i| = |S_{i-1}|+1, \\
\hat{c}_{\Pop} &= c_{\Pop} + |S_i| - |S_{i-1}| = 1 - 1 = 0 && \text{perché } |S_i| = |S_{i-1}|-1, \\
\hat{c}_{\MultiPop} &= c_{\MultiPop} + |S_i| - |S_{i-1}| = k' - k' = 0 && \text{perché } |S_i| = |S_{i-1}|-k',
\end{align*}
dove $k' = \min(|S_{i-1}|,k)$ è il numero di elementi effettivamente rimossi. Troviamo gli stessi costi del metodo degli accantonamenti (il potenziale $|S|$ è la somma dei crediti accantonati sugli elementi), e per il lemma~\ref{lem:potenziale}
\[
\sum_{i=1}^{n} c_i \le \sum_{i=1}^{n} \hat{c}_i = 2\cdot\#\Push \le 2n.
\]

\subsection{Contatore binario con \Increment}
Per il contatore, con $D_i$ configurazione dopo $i$ incrementi, definiamo il potenziale come il numero di bit uguali a $1$:
\[
\Phi(D_i) = \#\{\text{bit uguali a } 1 \text{ in } D_i\}.
\]
Se $D_0$ è la configurazione nulla si ha $\Phi(D_0) = 0$ e in generale $\Phi(D_i) \ge \Phi(D_0)$. Per calcolare il costo ammortizzato dell'$i$-esimo incremento osserviamo che
\[
c_i = \#\SetBit_i + \#\ResetBit_i, \qquad \Delta\Phi_i = \#\SetBit_i - \#\ResetBit_i
\]
dove $\#\SetBit_i$ e $\#\ResetBit_i$ sono il numero di bit portati a $1$ e a $0$ dall'$i$-esimo incremento, con $0 \le \#\SetBit_i \le 1$. Infatti ogni \SetBit{} aggiunge un bit uguale a $1$ e ogni \ResetBit{} ne toglie uno. Pertanto
\[
\hat{c}_i = c_i + \Delta\Phi_i = \big(\#\SetBit_i + \#\ResetBit_i\big) + \big(\#\SetBit_i - \#\ResetBit_i\big) = 2\cdot\#\SetBit_i \le 2
\]
da cui
\[
\sum_{i=1}^{n} c_i \le \sum_{i=1}^{n} \hat{c}_i = 2\cdot \sum_{i=1}^{n} \#\SetBit_i \le 2n.
\]

\subsection{Configurazione iniziale non nulla}
Il vantaggio del metodo del potenziale è che gestisce con facilità il caso in cui la struttura dati \emph{non} parta dalla configurazione iniziale ``vuota''. Dalla~\eqref{eq:telescopica} si ottiene infatti
\[
\sum_{i=1}^{n} c_i = \sum_{i=1}^{n} \hat{c}_i - \Phi(D_n) + \Phi(D_0)
\]
e basta maggiorare il potenziale iniziale.

\paragraph{Contatore non nullo.}
Usiamo la stessa funzione potenziale di prima. Si ha $\Phi(D_0) \le k$, dove $k$ è il numero di bit del contatore, ma in generale non vale più $\Phi(D_i) \ge \Phi(D_0)$ (il contatore può passare per configurazioni con meno bit uguali a $1$ di quelli iniziali). Il lemma~\ref{lem:potenziale} non è quindi applicabile, ma la~\eqref{eq:telescopica} dà comunque, poiché $\Phi(D_n) \ge 0$:
\[
\sum_{i=1}^{n} c_i = \sum_{i=1}^{n} \hat{c}_i - \Phi(D_n) + \Phi(D_0) \le 2n + \Phi(D_0) \le 2n + k.
\]
Se $k = O(n)$ allora $2n+k = O(n)$ e quindi $\sum c_i = O(n)$; in generale il costo è $O(n+k)$, come già trovato con l'aggregazione.

\paragraph{Stack non vuoto con \MultiPop.}
Con $\Phi(S) = |S|$ e uno stack iniziale $D_0$ che contiene $|D_0|$ elementi si ha $\Phi(D_0) = |D_0|$ e $\Phi(D_n) \ge 0$, quindi
\[
\sum_{i=1}^{n} c_i = \sum_{i=1}^{n} \hat{c}_i - \Phi(D_n) + \Phi(D_0) \le 2\cdot\#\Push + |D_0| \le 2n + |D_0|.
\]
Se $|D_0| = O(n)$, ovvero se lo stack iniziale non è troppo grande rispetto al numero di operazioni, allora $\sum c_i = O(n)$.

\begin{esercizio}[Costo pari a una potenza di $2$, con il potenziale]
Risolvere l'esercizio~\ref{es:potenze} con il metodo del potenziale.
\end{esercizio}

\paragraph{Soluzione.}
Poniamo $\Phi(D_0) = 0$ e, per $i \ge 1$,
\[
\Phi(D_i) = 2\left(i - 2^{\lfloor \lg i \rfloor}\right)
\]
che coincide con il credito $B_i$ accumulato dopo $i$ operazioni nella soluzione con gli accantonamenti, ed è sempre non negativa perché $2^{\lfloor \lg i \rfloor} \le i$. Calcoliamo i costi ammortizzati distinguendo i casi:
\begin{itemize}
    \item Se $i = 1$: $\hat{c}_1 = 1 + \Phi(D_1) - \Phi(D_0) = 1 + 0 - 0 = 1$.
    \item Se $i$ non è una potenza di $2$: $c_i = 1$ e $\lfloor \lg i \rfloor = \lfloor \lg (i-1) \rfloor$, quindi $\Phi(D_i) - \Phi(D_{i-1}) = 2$ e $\hat{c}_i = 1 + 2 = 3$.
    \item Se $i = 2^j$ con $j \ge 1$: $c_i = 2^j$, $\Phi(D_i) = 2(2^j - 2^j) = 0$ e $\Phi(D_{i-1}) = 2\left(2^j - 1 - 2^{j-1}\right) = 2^j - 2$, quindi $\hat{c}_i = 2^j + 0 - (2^j - 2) = 2$.
\end{itemize}
Si ritrovano gli stessi costi ammortizzati degli accantonamenti, e poiché $\Phi(D_n) \ge 0 = \Phi(D_0)$ vale $T(n) \le \sum \hat{c}_i \le 3n$: il costo ammortizzato per operazione è al più $3$.

\section{Confronto tra i metodi}\label{sec:confronto}
I tre metodi portano agli stessi risultati sugli esempi visti: $O(n)$ per $n$ operazioni su uno stack con \MultiPop{} e per $n$ incrementi di un contatore, quindi $O(1)$ ammortizzato per operazione. Cambia invece il modo di arrivarci, riassunto nella tabella seguente.

\begin{center}
\small
\begin{tabular}{p{2.6cm}p{3.5cm}p{3.5cm}p{3.5cm}}
\toprule
 & \textbf{Aggregazione} & \textbf{Accantonamenti} & \textbf{Potenziale} \\
\midrule
Costo ammortizzato & Uguale per tutte le operazioni: $T(n)/n$. & Scelto per ogni tipo di operazione. & $\hat{c}_i = c_i + \Delta\Phi_i$, può variare con l'operazione. \\[2mm]
Cosa si sceglie & Niente: si stima direttamente $T(n)$. & I costi $\hat{c}_i$. & La funzione $\Phi$. \\[2mm]
Dove sta il lavoro pagato in anticipo & Non viene tracciato. & Credito su singoli elementi. & Potenziale della struttura nel suo insieme. \\[2mm]
Cosa si verifica & Il calcolo di $T(n)$. & $B_j \ge 0$ per ogni $j$. & $\Phi(D_k) \ge \Phi(D_0)$ per ogni $k$. \\
\bottomrule
\end{tabular}
\end{center}

In sintesi, l'aggregazione è il metodo più semplice ma assegna lo stesso costo a tutte le operazioni, quindi funziona bene quando basta una stima del totale; gli accantonamenti permettono di distinguere le operazioni e hanno una lettura intuitiva in termini di monete, al prezzo di dover decidere dove si trova ogni unità di credito; il potenziale è il più generale, perché riassume tutto in una sola funzione della configurazione e gestisce con facilità il caso di una struttura dati che non parte da zero, ma richiede di trovare la funzione $\Phi$ giusta.

\section{Tabelle dinamiche}\label{sec:tabelle}
Molto spesso non si sa in anticipo quanti elementi dovrà contenere una tabella. Se si alloca poco spazio, prima o poi la tabella si riempie e un nuovo inserimento provoca un overflow; se se ne alloca troppo, gran parte della memoria resta inutilizzata. Le \textbf{tabelle dinamiche} risolvono il problema con la \textbf{riallocazione}: quando la tabella è piena se ne alloca una più grande, si copiano al suo interno tutti gli elementi e si libera quella vecchia; allo stesso modo, quando la tabella si svuota troppo, la si può sostituire con una più piccola. Una singola riallocazione è costosa, perché richiede di copiare tutti gli elementi, ma avviene raramente: è quindi la situazione ideale per l'analisi ammortizzata, che ci permetterà di dimostrare che inserimenti e cancellazioni hanno costo ammortizzato $O(1)$. Il modo in cui gli elementi sono organizzati all'interno della tabella (un array, uno heap, una tabella hash) non è importante ai fini dell'analisi.

\begin{definizione}[Fattore di carico]\label{def:carico}
Data una tabella $T$, indichiamo con $\tsize[T]$ la sua dimensione, ovvero il numero di posizioni disponibili, e con $\tnum[T]$ il numero di elementi che contiene. Il \textbf{fattore di carico} (load factor) di $T$ è
\[
\alpha(T) = \begin{cases} \dfrac{\tnum[T]}{\tsize[T]} & \text{se } \tsize[T] > 0 \\[3mm] 1 & \text{altrimenti.} \end{cases}
\]
\end{definizione}

Il fattore di carico misura quanto è piena la tabella: vale $1$ quando la tabella è piena e diminuisce man mano che si svuota. Per convenzione una tabella vuota, che ha dimensione $0$, è considerata piena. Se $\alpha(T)$ resta sempre al di sopra di una costante, lo spazio inutilizzato non supera mai una frazione costante dello spazio totale: è questa la garanzia che vogliamo ottenere, insieme a un costo ammortizzato costante per le operazioni.

\subsection{Inserimenti con espansione}\label{sec:tabelle-ins}
Consideriamo per ora soltanto gli inserimenti. Quando si inserisce un elemento in una tabella piena, la strategia più comune consiste nel \textbf{raddoppiare} la dimensione della tabella: si alloca una nuova tabella di dimensione $2\cdot\tsize[T]$, vi si copiano tutti gli elementi e si libera quella vecchia. Questa operazione è detta \textbf{espansione}, ed è descritta dall'algoritmo~\ref{alg:tableinsert}.

\begin{algorithm}[htbp]
\caption{L'operazione $\TableInsert(T,x)$ alloca la tabella al primo inserimento e la raddoppia (espansione) quando è piena, copiando tutti gli elementi nella nuova tabella prima di inserire $x$.}\label{alg:tableinsert}
\begin{algorithmic}[1]
\Procedure{Table-Insert}{$T,x$}
  \If{$\tsize[T] = 0$}
    \State allocare $\mathit{Table}[T]$ di dimensione $1$
    \State $\tsize[T] \gets 1$
    \State $\tnum[T] \gets 0$
  \ElsIf{$\tnum[T] = \tsize[T]$}
    \State allocare $\mathit{new}$ di dimensione $2\cdot\tsize[T]$
    \State copiare i $\tnum[T]$ elementi di $\mathit{Table}[T]$ in $\mathit{new}$\label{alg:tableinsert:copy}
    \State deallocare $\mathit{Table}[T]$
    \State $\mathit{Table}[T] \gets \mathit{new}$
    \State $\tsize[T] \gets 2\cdot\tsize[T]$
  \EndIf
  \State inserire $x$ in $\mathit{Table}[T]$\label{alg:tableinsert:ins}
  \State $\tnum[T] \gets \tnum[T] + 1$
\EndProcedure
\end{algorithmic}
\end{algorithm}

\paragraph{Il costo di un inserimento.}
Valutiamo i costi in termini di \textbf{inserimenti elementari}: costa $1$ l'inserimento di $x$ (riga~\ref{alg:tableinsert:ins}) e costa $1$ la copia di ciascun elemento durante un'espansione (riga~\ref{alg:tableinsert:copy}), mentre il costo di allocazioni e deallocazioni è costante o comunque dominato da quello delle copie. Un inserimento che non provoca un'espansione costa quindi $1$, mentre l'$i$-esimo inserimento, se provoca un'espansione, costa $i$: $i-1$ copie degli elementi già presenti più l'inserimento del nuovo elemento. Partendo da una tabella vuota, la dimensione è sempre una potenza di $2$, quindi la tabella è piena prima dell'$i$-esimo inserimento esattamente quando $i-1$ è una potenza di $2$:
\[
c_i = \begin{cases} i & \text{se } i-1 \in 2^{\mathbb{N}} \\ 1 & \text{altrimenti} \end{cases}
\]
dove $2^{\mathbb{N}} = \{1, 2, 4, 8, \dots\}$ è l'insieme delle potenze di $2$. La tabella seguente mostra i primi nove inserimenti.

\begin{center}
\small
\begin{tabular}{cccccc}
\toprule
$i$ & $\tnum$ prima & $\tsize$ prima & espansione & $\tsize$ dopo & $c_i$ \\
\midrule
$1$ & $0$ & $0$ & (allocazione) & $1$ & $1$ \\
$2$ & $1$ & $1$ & sì & $2$ & $2$ \\
$3$ & $2$ & $2$ & sì & $4$ & $3$ \\
$4$ & $3$ & $4$ & no & $4$ & $1$ \\
$5$ & $4$ & $4$ & sì & $8$ & $5$ \\
$6$ & $5$ & $8$ & no & $8$ & $1$ \\
$7$ & $6$ & $8$ & no & $8$ & $1$ \\
$8$ & $7$ & $8$ & no & $8$ & $1$ \\
$9$ & $8$ & $8$ & sì & $16$ & $9$ \\
\bottomrule
\end{tabular}
\end{center}

\paragraph{Analisi naive.}
Consideriamo $n = 2^k$ inserimenti a partire da una tabella vuota. L'inserimento di indice $2^{k-1}+1$ trova la tabella piena con $2^{k-1}$ elementi e costa
\[
2^{k-1} + 1 = \frac{n}{2} + 1,
\]
quindi il costo di un singolo inserimento nel caso peggiore è $O(n)$ e il costo di $n$ inserimenti è $n \cdot O(n) = O(n^2)$. Come negli esempi precedenti, la stima è troppo pessimistica: le espansioni sono rare, e tra un'espansione e la successiva il numero di inserimenti economici raddoppia.

\subsection{Analisi con l'aggregazione}
Il calcolo è molto simile a quello dell'esercizio~\ref{es:potenze}, con la differenza che le operazioni costose sono quelle di indice $i$ con $i-1$ potenza di $2$. Sia $n \ge 2$. Scriviamo il costo di ogni inserimento con espansione come $i = (i-1) + 1$, così da separare le copie dagli inserimenti veri e propri:
\[
T(n) = \sum_{\substack{i=1 \\ i-1 \in 2^{\mathbb{N}}}}^{n} i + \sum_{\substack{i=1 \\ i-1 \notin 2^{\mathbb{N}}}}^{n} 1 = \sum_{\substack{i=1 \\ i-1 \in 2^{\mathbb{N}}}}^{n} (i-1) + \sum_{\substack{i=1 \\ i-1 \in 2^{\mathbb{N}}}}^{n} 1 + \sum_{\substack{i=1 \\ i-1 \notin 2^{\mathbb{N}}}}^{n} 1.
\]
Le ultime due somme insieme contano tutte le $n$ operazioni, mentre nella prima i valori $i-1$ sono le potenze di $2$ comprese tra $1$ e $n-1$, ovvero $2^0, 2^1, \dots, 2^{\lfloor \lg(n-1) \rfloor}$. Si ottiene quindi
\[
T(n) = \sum_{j=0}^{\lfloor \lg(n-1) \rfloor} 2^j + n = 2^{\lfloor \lg(n-1) \rfloor + 1} - 1 + n \le 2(n-1) - 1 + n < 3n
\]
dove abbiamo usato $2^{\lfloor \lg(n-1) \rfloor} \le n-1$. Pertanto
\[
\text{costo ammortizzato per operazione} = \frac{T(n)}{n} < \frac{3n}{n} = 3 = O(1).
\]

\subsection{Analisi con gli accantonamenti}
Il risultato dell'aggregazione suggerisce di addebitare $3$ unità a ogni inserimento, e gli accantonamenti permettono di capire perché proprio $3$. Poniamo
\[
\hat{c}_{ins} = 3
\]
e ripartiamo le tre unità in questo modo:
\begin{itemize}
    \item \textbf{Inserimento}: $1$ unità paga il costo reale dell'inserimento elementare dell'elemento.
    \item \textbf{Copia dell'elemento}: $1$ unità viene accantonata sull'elemento stesso, e pagherà la sua copia alla prossima espansione.
    \item \textbf{Copia di un altro elemento}: $1$ unità viene accantonata su un elemento già ricopiato in un'espansione precedente, che ha speso il proprio credito e non avrebbe altrimenti come pagare la copia successiva.
\end{itemize}
Il primo inserimento non trova alcun elemento già ricopiato da aiutare, quindi per esso possiamo porre $\hat{c}_{ins} = 2$. Supponiamo che subito dopo un'espansione la tabella abbia dimensione $m$: contiene allora $m/2$ elementi, tutti senza credito, perché lo hanno appena speso per farsi copiare. Prima della prossima espansione verranno inseriti altri $m/2$ elementi, ognuno dei quali accantona $2$ unità, una su di sé e una su uno dei vecchi elementi: quando la tabella è di nuovo piena, ognuno degli $m$ elementi porta con sé esattamente l'unità di credito che serve a copiarlo. La figura~\ref{fig:crediti-tabella} mostra i primi sei inserimenti.

\begin{figure}[htbp]
\centering
\begin{widget}{dynamic-table}[steps=6]
\begin{tikzpicture}[
  cell/.style={draw=Accent, minimum width=6mm, minimum height=6mm, inner sep=0pt, font=\small},
  dk/.style={circle, fill=Accent, inner sep=0pt, minimum size=2.4mm},
  lt/.style={circle, fill=Accent!25, draw=Accent!60, inner sep=0pt, minimum size=2.4mm},
  lab/.style={font=\scriptsize, TextGray, anchor=north west}]
  % \celle{x}{y}{dimensione}{elementi}{etichetta}
  \newcommand{\celle}[5]{%
    \foreach \j in {1,...,#3}{\node[cell] at ({#1+0.6*(\j-1)},#2) {};}
    \foreach \e [count=\j] in {#4}{\node[font=\small] at ({#1+0.6*(\j-1)},#2) {$\e$};}
    \node[lab] at ({#1-0.3},{#2-0.4}) {#5};}
  % \cred{x}{y}{posizione}{stile}
  \newcommand{\cred}[4]{\node[#4] at ({#1+0.6*(#3-1)},{#2+0.55}) {};}
  % prima riga
  \celle{0}{0}{1}{a}{dopo $1$ ins.}       \cred{0}{0}{1}{dk}
  \celle{1.6}{0}{2}{a,b}{dopo $2$ ins.}   \cred{1.6}{0}{1}{lt} \cred{1.6}{0}{2}{dk}
  \celle{3.8}{0}{4}{a,b,c}{dopo $3$ ins.} \cred{3.8}{0}{2}{lt} \cred{3.8}{0}{3}{dk}
  \celle{7.2}{0}{4}{a,b,c,d}{dopo $4$ ins.} \cred{7.2}{0}{1}{lt} \cred{7.2}{0}{2}{lt} \cred{7.2}{0}{3}{dk} \cred{7.2}{0}{4}{dk}
  % seconda riga
  \celle{0}{-1.7}{8}{a,b,c,d,e}{dopo $5$ ins.}   \cred{0}{-1.7}{4}{lt} \cred{0}{-1.7}{5}{dk}
  \celle{5.6}{-1.7}{8}{a,b,c,d,e,f}{dopo $6$ ins.} \cred{5.6}{-1.7}{3}{lt} \cred{5.6}{-1.7}{4}{lt} \cred{5.6}{-1.7}{5}{dk} \cred{5.6}{-1.7}{6}{dk}
  % legenda
  \node[dk] at (0.4,-3.0) {}; \node[font=\scriptsize, anchor=west] at (0.6,-3.0) {credito su un elemento inserito e non ancora ricopiato};
  \node[lt] at (0.4,-3.5) {}; \node[font=\scriptsize, anchor=west] at (0.6,-3.5) {credito su un elemento già ricopiato};
\end{tikzpicture}
\end{widget}
\caption{Crediti accantonati durante i primi sei inserimenti, con $\hat{c}_{ins} = 3$ ($2$ per il primo). Ogni elemento inserito dopo l'ultima espansione porta una moneta scura e ne deposita una chiara su un elemento già ricopiato. Dopo $4$ inserimenti la tabella è piena e ognuno dei $4$ elementi ha una moneta: sono proprio quelle che pagano le $4$ copie dell'espansione provocata dal quinto inserimento, dopo la quale gli elementi $a,b,c,d$ restano senza credito finché i nuovi inserimenti non ne depositano altro.}
\label{fig:crediti-tabella}
\end{figure}

\paragraph{Il credito non è mai negativo.}
Verifichiamo ora formalmente che $B_j \ge 0$ per ogni $j$. Un inserimento senza espansione aumenta il credito di $\hat{c}_i - c_i = 3 - 1 = 2$, quindi l'unico momento in cui $B_j$ potrebbe diventare negativo è subito dopo un'espansione, ovvero dopo un inserimento di indice $2^k + 1$, che porta la tabella dalla dimensione $2^k$ alla dimensione $2^{k+1}$. Poiché $c_{2^k+1} = 2^k + 1$, si ha
\begin{equation}\label{eq:credito-espansione}
B_{2^k+1} = B_{2^k} + \hat{c}_{2^k+1} - c_{2^k+1} = B_{2^k} + 3 - (2^k + 1) = B_{2^k} - 2^k + 2
\end{equation}
quindi basta dimostrare che $B_{2^k} \ge 2^k - 2$ per ogni $k$. Supponiamo per assurdo che non sia così, e sia $m$ il più piccolo intero per cui
\[
B_{2^m} < 2^m - 2.
\]
Non può essere $m = 0$, perché $B_{2^0} = B_1 = \hat{c}_1 - c_1 = 2 - 1 = 1 > 2^0 - 2$; quindi $m > 0$ e, per la minimalità di $m$, vale $B_{2^{m-1}} \ge 2^{m-1} - 2$. Tra l'inserimento $2^{m-1}+1$ e l'inserimento $2^m$ ci sono $2^{m-1}-1$ inserimenti senza espansione, ognuno dei quali aggiunge $2$ al credito, quindi usando la~\eqref{eq:credito-espansione}
\[
\begin{aligned}
B_{2^m} &= B_{2^{m-1}+1} + 2\left(2^{m-1} - 1\right) = B_{2^{m-1}} - 2^{m-1} + 2 + 2^m - 2 \\
&= B_{2^{m-1}} + 2^{m-1} \ge 2^{m-1} - 2 + 2^{m-1} = 2^m - 2
\end{aligned}
\]
contraddicendo l'ipotesi. Pertanto $B_j \ge 0$ per ogni $j \ge 1$ e, per il lemma~\ref{lem:accantonamenti},
\[
T(n) = \sum_{i=1}^{n} c_i \le \sum_{i=1}^{n} \hat{c}_i \le 3n.
\]

\paragraph{Il credito in funzione di $\tnum$ e $\tsize$.}
Osservando la figura~\ref{fig:crediti-tabella} si può esprimere il credito accumulato nella tabella usando soltanto $\tnum[T]$ e $\tsize[T]$. Dopo l'ultima espansione sono stati inseriti $\tnum[T] - \tsize[T]/2$ elementi, ognuno dei quali ha depositato una moneta su di sé e una su un elemento già ricopiato, quindi il credito totale è
\[
B = 2\left(\tnum[T] - \frac{\tsize[T]}{2}\right) = 2\cdot\tnum[T] - \tsize[T].
\]
Ad esempio, dopo $6$ inserimenti si ha $\tnum = 6$ e $\tsize = 8$, e infatti in figura ci sono $2\cdot 6 - 8 = 4$ monete. La formula vale anche dopo il primo inserimento, che ha accantonato una sola unità: $2\cdot 1 - 1 = 1$. Questa espressione dipende soltanto dalla configurazione della tabella, ed è il punto di partenza del metodo del potenziale.

\subsection{Analisi con il potenziale}
Usiamo come funzione potenziale proprio il credito appena trovato:
\[
\Phi(T) = 2\cdot\tnum[T] - \tsize[T].
\]
Se $T_0$ è la tabella vuota si ha $\Phi(T_0) = 0$, e poiché la tabella è sempre piena almeno per metà, ovvero $\frac{1}{2}\tsize[T] \le \tnum[T]$, vale
\[
\Phi(T) \ge 0 = \Phi(T_0)
\]
per ogni configurazione $T$. Il metodo del potenziale può quindi essere utilizzato per stimare i costi reali. Il potenziale si comporta esattamente come serve: vale $0$ subito dopo un'espansione, quando $\tnum[T] = \tsize[T]/2$, e cresce di $2$ a ogni inserimento fino a valere $\tsize[T]$ quando la tabella è piena, cioè proprio il numero di elementi da copiare alla prossima espansione. La figura~\ref{fig:potenziale-tabella} mostra questo andamento.

\begin{figure}[htbp]
\centering
\begin{widget}{dynamic-table}[steps=32, coins=false, chart]
\begin{tikzpicture}[x=0.32cm, y=0.13cm]
  % griglia e assi
  \foreach \y in {8,16,24,32}{\draw[RuleGray, very thin] (0,\y) -- (32,\y); \node[font=\scriptsize, TextGray, anchor=east] at (-0.3,\y) {\y};}
  \foreach \x in {0,4,...,32}{\node[font=\scriptsize, TextGray, anchor=north] at (\x,-0.8) {\x};}
  \draw[->] (0,0) -- (33.5,0) node[right, font=\small] {$i$};
  \draw[->] (0,0) -- (0,35);
  % num, size, Phi
  \draw[Accent!60, thick] plot coordinates {(0,0) (32,32)};
  \draw[FigColor, thick, dashed] plot coordinates {(0,0) (1,1) (2,2) (3,4) (4,4) (5,8) (6,8) (7,8) (8,8) (9,16) (10,16) (11,16) (12,16) (13,16) (14,16) (15,16) (16,16) (17,32) (18,32) (19,32) (20,32) (21,32) (22,32) (23,32) (24,32) (25,32) (26,32) (27,32) (28,32) (29,32) (30,32) (31,32) (32,32)};
  \draw[LemmaColor, very thick] plot coordinates {(0,0) (1,1) (2,2) (3,2) (4,4) (5,2) (6,4) (7,6) (8,8) (9,2) (10,4) (11,6) (12,8) (13,10) (14,12) (15,14) (16,16) (17,2) (18,4) (19,6) (20,8) (21,10) (22,12) (23,14) (24,16) (25,18) (26,20) (27,22) (28,24) (29,26) (30,28) (31,30) (32,32)};
  % legenda
  \draw[Accent!60, thick] (35,24) -- (37,24) node[right, font=\small, black] {$\tnum_i$};
  \draw[FigColor, thick, dashed] (35,19) -- (37,19) node[right, font=\small, black] {$\tsize_i$};
  \draw[LemmaColor, very thick] (35,14) -- (37,14) node[right, font=\small, black] {$\Phi_i$};
\end{tikzpicture}
\end{widget}
\caption{Andamento di $\tnum_i$, $\tsize_i$ e $\Phi_i = 2\cdot\tnum_i - \tsize_i$ durante $32$ inserimenti in una tabella inizialmente vuota, misurati dopo l'$i$-esimo inserimento. Subito prima di ogni espansione il potenziale ha raggiunto il numero di elementi presenti nella tabella, e può quindi pagare la loro copia; subito dopo scende a $0$, ma l'inserimento che ha provocato l'espansione lo riporta a $2$.}
\label{fig:potenziale-tabella}
\end{figure}

Indichiamo con $\tnum_i$ e $\tsize_i$ il numero di elementi e la dimensione della tabella dopo l'$i$-esimo inserimento, e con $\Phi_i = 2\cdot\tnum_i - \tsize_i$ il potenziale corrispondente. In ogni caso $\tnum_i = \tnum_{i-1} + 1$, e distinguiamo tre casi.

\paragraph{Inserimento senza espansione.}
La dimensione non cambia, ovvero $\tsize_i = \tsize_{i-1}$, e il costo reale è $1$:
\[
\begin{aligned}
\hat{c}_{ins} &= 1 + (2\cdot\tnum_i - \tsize_i) - (2\cdot\tnum_{i-1} - \tsize_{i-1}) \\
&= 1 + 2(\tnum_{i-1} + 1) - \tsize_{i-1} - 2\cdot\tnum_{i-1} + \tsize_{i-1} = 3.
\end{aligned}
\]

\paragraph{Inserimento con espansione in una tabella non vuota.}
La tabella era piena, quindi $\tsize_{i-1} = \tnum_{i-1}$, e dopo l'espansione $\tsize_i = 2\cdot\tsize_{i-1} = 2\cdot\tnum_{i-1}$. Il costo reale è $\tnum_i = \tnum_{i-1} + 1$, ovvero $\tnum_{i-1}$ copie più l'inserimento:
\[
\begin{aligned}
\hat{c}_{ins} &= (\tnum_{i-1} + 1) + \bigl(2(\tnum_{i-1} + 1) - 2\cdot\tnum_{i-1}\bigr) - (2\cdot\tnum_{i-1} - \tnum_{i-1}) \\
&= \tnum_{i-1} + 1 + 2 - \tnum_{i-1} = 3.
\end{aligned}
\]

\paragraph{Inserimento in una tabella vuota.}
È il primo inserimento ($i = 1$): si ha $\tsize_1 = \tnum_1 = 1$ e $\Phi_0 = 0$, quindi
\[
\hat{c}_{ins} = 1 + (2\cdot 1 - 1) - 0 = 2.
\]

Ritroviamo così gli stessi costi degli accantonamenti, e poiché $\Phi_i \ge \Phi_0$
\[
T(n) = \sum_{i=1}^{n} c_i \le \sum_{i=1}^{n} \hat{c}_i \le 3n.
\]
Notiamo infine che, poiché la tabella è sempre piena almeno per metà, il fattore di carico non scende mai sotto $1/2$: lo spazio sprecato non supera mai la metà dello spazio allocato.

\subsection{Inserimenti e cancellazioni}
Per non sprecare memoria, quando si cancellano molti elementi conviene \textbf{contrarre} la tabella: si alloca una tabella più piccola, vi si copiano gli elementi rimasti e si libera quella vecchia. L'operazione $\TableDelete(T,x)$ è analoga a \TableInsert: rimuove l'elemento, con costo $1$, ed esegue una \textbf{contrazione} quando il fattore di carico diventa troppo basso. Una contrazione costa quanto il numero di elementi da copiare, quindi una cancellazione che la provoca ha costo $\tnum_{i-1}$, ovvero $\tnum_{i-1} - 1$ copie più la cancellazione stessa.

\paragraph{Dimezzare a $1/2$ non funziona.}
La strategia più naturale sembra la simmetrica dell'espansione: raddoppiare la tabella quando è piena e dimezzarla quando il fattore di carico scende sotto $1/2$. In questo modo $\alpha(T)$ non scende mai sotto $1/2$, ma il costo ammortizzato può diventare molto alto. Consideriamo $n = 2^k$ operazioni: le prime $n/2$ sono inserimenti (I), che riempiono esattamente una tabella di dimensione $n/2$, e le restanti $n/2$ seguono lo schema I, C, C, I, I, C, C, I, I, \dots, dove C indica una cancellazione. La tabella seguente mostra cosa succede.

\begin{center}
\small
\begin{tabular}{lccccc}
\toprule
 & \multicolumn{1}{c}{I \dots I} & I & C & C & I \\
\midrule
$\tnum[T]$ & $n/2$ & $n/2 + 1$ & $n/2$ & $n/2 - 1$ & $n/2$ \\
$\tsize[T]$ & $n/2$ & $n$ & $n$ & $n/2$ & $n/2$ \\
$\alpha(T)$ & $1$ & $> 1/2$ & $= 1/2$ & $< 1$ & $1$ \\
$c_i$ & $< 3n/2$ in totale & $n/2 + 1$ & $1$ & $n/2$ & $1$ \\
\bottomrule
\end{tabular}
\end{center}

\noindent
Il primo inserimento del secondo gruppo provoca un'espansione, la prima cancellazione riporta il fattore di carico esattamente a $1/2$ e la seconda lo fa scendere sotto, provocando una contrazione; dopo il successivo inserimento la tabella è di nuovo piena, e l'inserimento seguente provoca un'altra espansione. Ogni gruppo di quattro operazioni contiene quindi un'espansione e una contrazione, che costano circa $n/2$ ciascuna, e i gruppi sono $n/8$: il costo totale è dell'ordine di $\frac{n}{8}\cdot n = \frac{n^2}{8}$, ovvero $\Theta(n^2)$, e il costo ammortizzato per operazione è $\Theta(n)$. Il problema è che subito dopo un'espansione non si cancellano abbastanza elementi da pagare una contrazione, e subito dopo una contrazione non se ne inseriscono abbastanza da pagare un'espansione.

\paragraph{Dimezzare a $1/4$.}
La soluzione è lasciare più margine: si continua a raddoppiare la tabella quando è piena, ma la si dimezza soltanto quando una cancellazione farebbe scendere il fattore di carico sotto $1/4$. In altre parole ci accontentiamo di avere
\[
\alpha(T) \ge \frac{1}{4}
\]
in cambio di un costo ammortizzato costante.\footnote{Per semplicità trascuriamo il caso limite in cui la tabella si svuota del tutto; nel libro si assume che in quel caso la tabella venga deallocata, ponendo $\tsize[T] = 0$.} Sia dopo un'espansione sia dopo una contrazione il fattore di carico torna circa a $1/2$, e serve un numero di operazioni proporzionale a $\tnum[T]$ per arrivare alla riallocazione successiva, in un senso o nell'altro.

Ci serve quindi un potenziale che valga $0$ quando $\alpha(T) = 1/2$ e cresca man mano che il fattore di carico si allontana da $1/2$, in modo da avere accumulato abbastanza potenziale da pagare la copia di tutti gli elementi quando $\alpha(T)$ arriva a $1$ oppure scende a $1/4$. Poniamo
\[
\Phi(T) = \begin{cases} 2\cdot\tnum[T] - \tsize[T] & \text{se } \alpha(T) \ge \dfrac{1}{2} \\[3mm] \dfrac{1}{2}\tsize[T] - \tnum[T] & \text{se } \alpha(T) \le \dfrac{1}{2}. \end{cases}
\]
Le due espressioni coincidono quando $\alpha(T) = 1/2$: in quel caso $\tnum[T] = \frac{1}{2}\tsize[T]$, quindi $\frac{1}{2}\tsize[T] - \tnum[T] = 0$ e $2\cdot\tnum[T] - \tsize[T] = 0$. Questo ci permette di usare l'una o l'altra, a seconda di cosa è più comodo, quando la tabella è piena esattamente per metà. Inoltre vale sempre $\Phi(T) \ge 0 = \Phi(T_0)$, perché ciascuna delle due espressioni è non negativa nel proprio intervallo. La figura~\ref{fig:ricarica} mostra l'andamento del potenziale in funzione del fattore di carico.

\begin{figure}[htbp]
\centering
\begin{widget}{table-potential}
\begin{tikzpicture}[x=9cm, y=3.2cm]
  % zone
  \fill[ObsColor!8] (0.25,0) rectangle (0.5,1.1);
  \fill[Accent!6] (0.5,0) rectangle (1,1.1);
  % assi
  \draw[->] (0,0) -- (1.08,0) node[right, font=\small] {$\alpha(T)$};
  \draw[->] (0,0) -- (0,1.2) node[above, font=\small] {$\Phi(T)/\tsize[T]$};
  \foreach \x/\l in {0.25/\frac14, 0.5/\frac12, 1/1}{\draw (\x,0.02) -- (\x,-0.02) node[below, font=\scriptsize] {$\l$};}
  \foreach \y/\l in {0.25/\frac14, 1/1}{\draw (0.01,\y) -- (-0.01,\y) node[left, font=\scriptsize] {$\l$};}
  % Phi = num (riferimento)
  \draw[TextGray, dashed] (0,0) -- (1.05,1.05);
  \node[font=\scriptsize, TextGray, rotate=19.6] at (0.80,0.87) {$\Phi = \tnum[T]$};
  % potenziale
  \draw[LemmaColor, very thick] (0.25,0.25) -- (0.5,0) -- (1,1);
  \fill[LemmaColor] (0.25,0.25) circle (0.6pt);
  \fill[LemmaColor] (1,1) circle (0.6pt);
  % etichette delle zone
  \node[font=\scriptsize, align=center, ObsColor] at (0.375,0.6) {il potenziale\\paga la\\contrazione};
  \node[font=\scriptsize, align=center, Accent] at (0.82,0.22) {il potenziale\\paga l'espansione};
\end{tikzpicture}
\end{widget}
\caption{Potenziale della tabella, in rapporto alla sua dimensione, in funzione del fattore di carico. Il potenziale è nullo quando la tabella è piena per metà e cresce linearmente in entrambe le direzioni: quando $\alpha(T) = 1$ il potenziale vale $\tsize[T]$, cioè $\tnum[T]$, e quando $\alpha(T) = 1/4$ vale $\frac{1}{4}\tsize[T]$, cioè ancora $\tnum[T]$. In entrambi i casi tocca la retta tratteggiata $\Phi = \tnum[T]$, cioè è sufficiente a pagare la copia di tutti gli elementi.}
\label{fig:ricarica}
\end{figure}

Indichiamo come prima con $\tnum_i$ e $\tsize_i$ i valori dopo l'$i$-esima operazione e con $\alpha_{i-1} = \alpha(T_{i-1})$ il fattore di carico prima dell'operazione, e calcoliamo i costi ammortizzati distinguendo i casi.

\paragraph{Analisi degli inserimenti.}
In ogni caso $\tnum_i = \tnum_{i-1} + 1$.
\begin{itemize}
    \item \textbf{Caso $\alpha_{i-1} \ge 1/2$}: anche dopo l'inserimento la tabella è piena almeno per metà, quindi si usa la prima espressione del potenziale sia prima sia dopo. Il calcolo è lo stesso già fatto nel caso dei soli inserimenti, e dà $\hat{c}_{ins} = 3$ (oppure $2$ per il primo inserimento).
    \item \textbf{Caso $1/4 \le \alpha_{i-1} < 1/2$}: non c'è espansione, quindi $\tsize_i = \tsize_{i-1}$, e poiché $\tnum_{i-1} < \frac{1}{2}\tsize_{i-1}$ si ha $\tnum_i \le \frac{1}{2}\tsize_i$, ovvero $\alpha_i \le 1/2$. Si usa allora la seconda espressione:
    \[
    \begin{aligned}
    \hat{c}_{ins} &= 1 + \left(\tfrac{1}{2}\tsize_i - \tnum_i\right) - \left(\tfrac{1}{2}\tsize_{i-1} - \tnum_{i-1}\right) \\
    &= 1 + \tfrac{1}{2}\tsize_{i-1} - \tnum_{i-1} - 1 - \tfrac{1}{2}\tsize_{i-1} + \tnum_{i-1} = 0.
    \end{aligned}
    \]
\end{itemize}
Un inserimento in una tabella piena per meno di metà ha costo ammortizzato nullo: il suo costo reale è pagato dalla diminuzione del potenziale, perché avvicina il fattore di carico a $1/2$.

\paragraph{Analisi delle cancellazioni.}
In ogni caso $\tnum_i = \tnum_{i-1} - 1$.
\begin{itemize}
    \item \textbf{Caso $1/2 < \alpha_{i-1} \le 1$}: non c'è contrazione e dopo la cancellazione la tabella è ancora piena almeno per metà, quindi si usa la prima espressione con $\tsize_i = \tsize_{i-1}$:
    \[
    \begin{aligned}
    \hat{c}_{canc} &= 1 + (2\cdot\tnum_i - \tsize_i) - (2\cdot\tnum_{i-1} - \tsize_{i-1}) \\
    &= 1 + 2\cdot\tnum_{i-1} - 2 - \tsize_{i-1} - 2\cdot\tnum_{i-1} + \tsize_{i-1} = -1.
    \end{aligned}
    \]
    \item \textbf{Caso $1/4 < \alpha_{i-1} \le 1/2$}: non c'è contrazione e dopo la cancellazione $\alpha_i < 1/2$, quindi si usa la seconda espressione con $\tsize_i = \tsize_{i-1}$:
    \[
    \begin{aligned}
    \hat{c}_{canc} &= 1 + \left(\tfrac{1}{2}\tsize_i - \tnum_i\right) - \left(\tfrac{1}{2}\tsize_{i-1} - \tnum_{i-1}\right) \\
    &= 1 + \tfrac{1}{2}\tsize_{i-1} - \tnum_{i-1} + 1 - \tfrac{1}{2}\tsize_{i-1} + \tnum_{i-1} = 2.
    \end{aligned}
    \]
    \item \textbf{Caso $\alpha_{i-1} = 1/4$}: la cancellazione provoca una contrazione. Si ha $\tnum_{i-1} = \frac{1}{4}\tsize_{i-1}$ e $\tsize_i = \frac{1}{2}\tsize_{i-1}$, e il costo reale è $\tnum_{i-1}$. Anche dopo la contrazione $\alpha_i < 1/2$, quindi
    \[
    \begin{aligned}
    \hat{c}_{canc} &= \tnum_{i-1} + \left(\tfrac{1}{2}\tsize_i - \tnum_i\right) - \left(\tfrac{1}{2}\tsize_{i-1} - \tnum_{i-1}\right) \\
    &= \tfrac{1}{4}\tsize_{i-1} + \left(\tfrac{1}{4}\tsize_{i-1} - \tfrac{1}{4}\tsize_{i-1} + 1\right) - \left(\tfrac{1}{2}\tsize_{i-1} - \tfrac{1}{4}\tsize_{i-1}\right) \\
    &= \tfrac{1}{4}\tsize_{i-1} + 1 - \tfrac{1}{4}\tsize_{i-1} = 1.
    \end{aligned}
    \]
\end{itemize}

\paragraph{Riassumendo.}
I costi ammortizzati delle due operazioni sono
\[
\hat{c}_{ins} = \begin{cases} 0 & \text{se } \frac{1}{4} \le \alpha(T) < \frac{1}{2} \\[1mm] 3 & \text{se } \frac{1}{2} \le \alpha(T) \le 1 \end{cases}
\qquad\qquad
\hat{c}_{canc} = \begin{cases} -1 & \text{se } \frac{1}{2} < \alpha(T) \le 1 \\[1mm] 2 & \text{se } \frac{1}{4} < \alpha(T) \le \frac{1}{2} \\[1mm] 1 & \text{se } \alpha(T) = \frac{1}{4} \end{cases}
\]
dove $\alpha(T)$ è il fattore di carico prima dell'operazione. Pertanto $\hat{c}_{ins} \le 3$ e $\hat{c}_{canc} \le 2 < 3$, e poiché $\Phi(T_i) \ge 0 = \Phi(T_0)$ si ottiene
\[
T(n) = \sum_{i=1}^{n} c_i \le \sum_{i=1}^{n} \hat{c}_i \le 3n.
\]
Qualsiasi sequenza di $n$ inserimenti e cancellazioni su una tabella inizialmente vuota costa quindi $O(n)$, e il costo ammortizzato di ciascuna operazione è $O(1)$.

In sintesi, raddoppiare la tabella quando è piena garantisce un costo ammortizzato costante per gli inserimenti e uno spreco di memoria al più pari a metà dello spazio allocato; dimezzarla appena scende sotto metà sembra la scelta simmetrica, ma porta a sequenze di costo $\Theta(n^2)$; dimezzarla soltanto sotto $1/4$ mantiene il costo ammortizzato costante anche con le cancellazioni, al prezzo di uno spreco di memoria che può arrivare a tre quarti dello spazio allocato.

\section*{Riferimenti}
Il materiale visto in questo capitolo comprende:
\begin{itemize}
    \item Materiale visto a lezione.
    \item Capitolo 17 (\emph{Amortized Analysis}, sezioni 17.1--17.4) del libro T.\,H.~Cormen, C.\,E.~Leiserson, R.\,L.~Rivest, C.~Stein, \emph{Introduction to Algorithms}.
\end{itemize}
